% ============================================================
%  CHAPITRE 2 — ÉTAT DE L'ART
% ============================================================
\chapter{État de l'art}

\section{Introduction}

Ce chapitre présente les fondements théoriques sur lesquels repose notre système de
maintenance prédictive. Nous abordons successivement les stratégies de maintenance,
la technologie SMART, les algorithmes d'apprentissage automatique utilisés (Random
Forest, Isolation Forest, LSTM), les modèles de langage de grande taille (LLM), et
les métriques d'évaluation des systèmes conversationnels.

% ─────────────────────────────────────────────────────────────
\section{Maintenance prédictive — Contexte et enjeux}
% ─────────────────────────────────────────────────────────────

\subsection{Évolution des stratégies de maintenance}

La maintenance des équipements informatiques a évolué à travers trois générations
\cite{mobley2002maintenance} :

\begin{description}
  \item[\textbf{Maintenance corrective (réactive)}] : on intervient uniquement après
  la panne. Cette approche est coûteuse en temps d'arrêt et en ressources humaines,
  et ne permet aucune anticipation.

  \item[\textbf{Maintenance préventive (planifiée)}] : les interventions sont
  programmées à intervalles fixes, indépendamment de l'état réel de l'équipement.
  Cela peut entraîner des remplacements prématurés ou des pannes entre deux
  interventions planifiées.

  \item[\textbf{Maintenance prédictive (conditionnelle)}] : l'état réel de
  l'équipement est surveillé en continu. On n'intervient que lorsque des signes de
  dégradation sont détectés. Cette approche optimise les coûts et réduit les temps
  d'arrêt non planifiés.
\end{description}

\subsection{Données SMART — Self-Monitoring, Analysis and Reporting Technology}

La technologie SMART est intégrée dans la quasi-totalité des disques durs modernes
depuis les années 1990. Elle permet de surveiller en temps réel l'état de santé du
disque à travers une série d'attributs normalisés. Dans notre système, le module
\texttt{SmartReader} utilise \texttt{smartctl} pour lire ces attributs sur Windows
et Linux.

\begin{table}[H]
\centering
\caption{Principaux attributs SMART utilisés dans le projet}
\begin{tabularx}{\textwidth}{|c|l|X|}
\hline
\rowcolor{lightblue}
\textbf{ID} & \textbf{Nom} & \textbf{Rôle dans notre système} \\
\hline
5   & Reallocated Sectors Count & Mappé sur \texttt{read\_errors} \\
\hline
187 & Reported Uncorrectable Errors & Mappé sur \texttt{write\_errors} \\
\hline
194 & Temperature Celsius & Mappé sur \texttt{temperature} \\
\hline
197 & Current Pending Sector Count & Contribue au \texttt{health\_score} LSTM \\
\hline
198 & Offline Uncorrectable & Contribue au \texttt{health\_score} LSTM \\
\hline
\end{tabularx}
\end{table}

Des études ont montré que les attributs 5, 187, 197 et 198 sont les meilleurs
prédicteurs de défaillance imminente d'un disque dur \cite{backblaze2020,
pinheiro2007smart}. Notre modèle LSTM exploite précisément ces attributs, mappés
depuis le dataset Backblaze vers notre schéma de base de données.

Lorsque \texttt{smartctl} n'est pas disponible (machines virtuelles, conteneurs),
l'agent utilise des valeurs de fallback : \texttt{health\_status = GOOD},
\texttt{read\_errors = 0}, \texttt{write\_errors = 0}, \texttt{temperature = 40°C}.
Cela garantit que la collecte continue même sur des environnements sans accès SMART.

% ─────────────────────────────────────────────────────────────
\section{Apprentissage automatique pour la maintenance prédictive}
% ─────────────────────────────────────────────────────────────

\subsection{Formulation du problème}

La prédiction de panne est traitée comme un problème de \textbf{classification
binaire} : à partir des métriques historiques d'une machine sur les 30 derniers
jours, le modèle prédit si une panne est probable dans un horizon de 7, 14 ou
30 jours. La sortie est une probabilité entre 0 et 100\%, convertie ensuite en
niveau de risque (LOW / MEDIUM / HIGH / CRITICAL).

\subsection{Random Forest}

Le Random Forest \cite{breiman2001rf} est un algorithme d'ensemble qui combine
plusieurs arbres de décision entraînés sur des sous-ensembles aléatoires des
données. La prédiction finale est obtenue par vote majoritaire entre tous les
arbres. Plus un arbre vote pour la classe « panne », plus la probabilité estimée
est élevée.

\textbf{Pourquoi ce choix pour notre projet :}
\begin{itemize}
  \item Robuste aux valeurs aberrantes et aux données manquantes ;
  \item Fournit un score d'importance pour chaque feature, affiché dans le
        composant \texttt{MachineDetails} du frontend ;
  \item Performant sur les données tabulaires sans nécessiter de GPU ;
  \item Gère nativement les classes déséquilibrées via le paramètre
        \texttt{class\_weight='balanced'} ;
  \item Facile à versionner et à comparer entre cycles d'entraînement.
\end{itemize}

\textbf{Hyperparamètres retenus :}
\begin{table}[H]
\centering
\caption{Hyperparamètres du Random Forest}
\begin{tabularx}{\textwidth}{|l|c|X|}
\hline
\rowcolor{lightblue}
\textbf{Paramètre} & \textbf{Valeur} & \textbf{Justification} \\
\hline
\texttt{n\_estimators} & 100 & Bon compromis précision / temps de calcul \\
\hline
\texttt{max\_depth} & 10 & Évite l'overfitting sur 63 features \\
\hline
\texttt{min\_samples\_split} & 5 & Régularisation des nœuds internes \\
\hline
\texttt{class\_weight} & balanced & Compense la rareté des pannes dans les données \\
\hline
\texttt{random\_state} & 42 & Reproductibilité des résultats \\
\hline
\end{tabularx}
\end{table}

\subsection{Isolation Forest pour la détection d'anomalies}

L'Isolation Forest \cite{liu2008isolation} est un algorithme non supervisé de
détection d'anomalies. Son principe est simple : les observations anormales sont
rares et différentes des observations normales, donc elles sont plus faciles à
isoler. L'algorithme construit des arbres aléatoires et mesure combien d'étapes
sont nécessaires pour isoler chaque point — moins il en faut, plus le point est
anormal.

Dans notre système, l'Isolation Forest est configuré avec 10\% d'anomalies
attendues (\texttt{contamination = 0.1}). Il opère sur les mêmes 63 features que
le Random Forest, mais de manière non supervisée, ce qui permet de détecter des
comportements inhabituels même sans historique de pannes.

\subsection{Ingénierie des caractéristiques (Feature Engineering)}

L'extraction de caractéristiques pertinentes est une étape cruciale du pipeline ML.
La classe \texttt{FeatureExtractor} extrait \textbf{63 caractéristiques} par machine
à partir des 30 derniers jours de données stockées dans PostgreSQL.

\begin{table}[H]
\centering
\caption{Les 63 caractéristiques extraites par \texttt{FeatureExtractor}}
\begin{tabularx}{\textwidth}{|c|l|X|}
\hline
\rowcolor{lightblue}
\textbf{Nb} & \textbf{Type} & \textbf{Description} \\
\hline
45 & Statistiques glissantes & Moyenne, médiane, écart-type, min, max sur 3 métriques
     $\times$ 3 fenêtres (24h, 168h, 720h) \\
\hline
3  & Tendances & Pente de la droite de régression sur 7 jours (via \texttt{numpy.polyfit}) \\
\hline
3  & Taux de variation & Différence entre la valeur courante et la valeur de l'heure précédente \\
\hline
3  & Volatilité & Coefficient de variation (écart-type / moyenne) \\
\hline
6  & SMART & Température, erreurs lecture/écriture, statut de santé (encodage one-hot) \\
\hline
3  & Temporel & Heure de la journée, jour de la semaine, indicateur week-end \\
\hline
\end{tabularx}
\end{table}

Les fenêtres temporelles de 24h, 168h (7 jours) et 720h (30 jours) permettent au
modèle de capturer à la fois les variations à court terme et les tendances de
dégradation à long terme. Les valeurs manquantes sont traitées par forward-fill
(propagation de la dernière valeur connue, limitée à 3 points consécutifs), puis
remplacement par zéro.

% ─────────────────────────────────────────────────────────────
\section{Réseaux de neurones récurrents et LSTM}
% ─────────────────────────────────────────────────────────────

\subsection{Réseaux de neurones récurrents (RNN)}

Les RNN sont conçus pour traiter des séquences de données en conservant une
mémoire des pas de temps précédents. Contrairement aux réseaux classiques, ils
maintiennent un état interne qui se met à jour à chaque nouvelle entrée. Cependant,
les RNN classiques ont du mal à mémoriser des informations sur de longues séquences
— c'est le problème de la disparition du gradient.

\subsection{Long Short-Term Memory (LSTM)}

Le LSTM \cite{hochreiter1997lstm} résout ce problème grâce à un mécanisme de
\textbf{portes} qui contrôle sélectivement ce que le réseau doit mémoriser, oublier
ou transmettre :

\begin{itemize}
  \item \textbf{Porte d'oubli} : décide quelle information de l'état précédent
        doit être effacée ;
  \item \textbf{Porte d'entrée} : décide quelle nouvelle information doit être
        mémorisée ;
  \item \textbf{Porte de sortie} : décide quelle partie de la mémoire doit être
        transmise comme sortie.
\end{itemize}

Ce mécanisme permet au LSTM de capturer des dépendances à long terme dans les
séquences de métriques SMART, comme une augmentation progressive des erreurs sur
plusieurs jours précédant une panne.

\subsection{Architecture LSTM dans notre projet}

Notre modèle PyTorch (\texttt{lstm\_predictor.py}) analyse des séquences de
\textbf{5 mesures SMART consécutives} pour prédire la probabilité de défaillance
d'un disque. L'architecture est volontairement simple et légère :

\begin{itemize}
  \item \textbf{Entrée} : séquence de 5 pas de temps, chacun avec 4 valeurs :
        \texttt{read\_errors}, \texttt{write\_errors}, \texttt{temperature},
        \texttt{health\_score} ;
  \item \textbf{Couche LSTM} : 32 unités cachées ;
  \item \textbf{Couche de sortie} : une valeur entre 0 et 1 (probabilité de panne),
        obtenue via une activation sigmoïde.
\end{itemize}

Le \texttt{health\_score} est un encodage numérique du statut SMART :
0,0 pour GOOD, 0,5 pour WARNING, 1,0 pour CRITICAL. Les quatre features sont
normalisées par des valeurs maximales fixes avant d'être passées au modèle, ce
qui évite les sorties saturées.

\subsection{Dataset Backblaze}

Le dataset Backblaze est une référence dans le domaine de la prédiction de
défaillance des disques durs. Backblaze publie trimestriellement les données SMART
de ses disques en production, avec les événements de panne réels documentés.

Pour le Q1 2020, le dataset utilisé dans notre projet contient :
\begin{itemize}
  \item 874 892 lignes de données SMART réelles ;
  \item 4 993 disques uniques ;
  \item 9 981 séquences temporelles construites (fenêtre glissante de 5 jours) ;
  \item Des pannes réelles documentées (colonne \texttt{failure} = 1).
\end{itemize}

\textbf{Correspondance des attributs Backblaze vers notre schéma :}

\begin{table}[H]
\centering
\caption{Mapping des attributs Backblaze vers les features LSTM}
\begin{tabular}{|l|l|l|}
\hline
\rowcolor{lightblue}
\textbf{Attribut Backblaze} & \textbf{Signification SMART} & \textbf{Feature projet} \\
\hline
\texttt{smart\_5\_raw} & Reallocated Sectors Count & \texttt{read\_errors} \\
\hline
\texttt{smart\_187\_raw} & Reported Uncorrectable Errors & \texttt{write\_errors} \\
\hline
\texttt{smart\_194\_raw} & Temperature Celsius & \texttt{temperature} \\
\hline
\texttt{smart\_197\_raw} + \texttt{smart\_198\_raw} & Pending + Uncorrectable Sectors & \texttt{health\_score} \\
\hline
\end{tabular}
\end{table}

\subsection{Gestion du déséquilibre de classes}

Le dataset Backblaze est fortement déséquilibré : les pannes représentent moins
de 1\% des enregistrements. Trois techniques sont combinées pour gérer ce problème :

\begin{enumerate}
  \item \textbf{Oversampling} : les séquences de panne sont dupliquées jusqu'à
        représenter 10\% du dataset d'entraînement ;
  \item \textbf{Pondération de la loss} : la fonction de perte pénalise 50 fois
        plus fortement les pannes manquées que les fausses alarmes ;
  \item \textbf{Réduction progressive du learning rate} : le taux d'apprentissage
        est divisé par 2 tous les 7 epochs pour stabiliser la convergence.
\end{enumerate}

% ─────────────────────────────────────────────────────────────
\section{Modèles de langage de grande taille (LLM)}
% ─────────────────────────────────────────────────────────────

\subsection{Architecture Transformer}

Les LLM modernes reposent sur l'architecture Transformer \cite{vaswani2017attention},
qui utilise un mécanisme d'\textbf{attention} pour pondérer l'importance de chaque
mot dans une phrase par rapport aux autres. Contrairement aux RNN, le Transformer
traite tous les mots en parallèle, ce qui le rend beaucoup plus rapide à entraîner
et capable de capturer des relations à très longue portée dans le texte.

\subsection{Ollama et LLaMA 3.2:1b}

Ollama est un framework permettant d'exécuter des LLM localement, sans connexion
internet ni GPU dédié \cite{touvron2023llama}. Dans notre projet, nous utilisons
\textbf{LLaMA 3.2:1b} (1 milliard de paramètres), qui offre un bon compromis entre
qualité des réponses et vitesse d'inférence sur du matériel standard.

Le choix d'un LLM local est motivé par la \textbf{confidentialité} : aucune donnée
sur l'infrastructure (noms de machines, métriques, alertes) n'est transmise à des
services cloud tiers.

\textbf{Paramètres d'inférence configurés :}
\begin{itemize}
  \item \texttt{temperature = 0.3} : réponses cohérentes et reproductibles,
        évitant les hallucinations sur des données techniques ;
  \item \texttt{num\_predict = 150} : limite la longueur des réponses pour
        la rapidité ;
  \item \texttt{repeat\_penalty = 1.1} : évite les répétitions dans les réponses.
\end{itemize}

% ─────────────────────────────────────────────────────────────
\section{Métriques d'évaluation}
% ─────────────────────────────────────────────────────────────

\subsection{ROUGE pour l'évaluation du chatbot}

ROUGE (\textit{Recall-Oriented Understudy for Gisting Evaluation})
\cite{lin2004rouge} est une famille de métriques qui mesure le chevauchement
lexical entre une réponse générée et une réponse de référence. Plus les mots
en commun sont nombreux, plus le score est élevé.

Notre pipeline d'évaluation (\texttt{evaluate.py}) calcule deux variantes :
\begin{itemize}
  \item \textbf{ROUGE-1} : chevauchement des mots individuels (unigrammes) ;
  \item \textbf{ROUGE-2} : chevauchement des paires de mots consécutifs (bigrammes).
\end{itemize}

Le score F1 est utilisé, qui équilibre précision (mots générés corrects) et rappel
(mots de référence retrouvés). La bibliothèque \texttt{rouge-score} est utilisée
avec \texttt{use\_stemmer=True} pour normaliser les formes fléchies en français.

\subsection{Métriques de classification pour les modèles ML}

\begin{table}[H]
\centering
\caption{Métriques d'évaluation des modèles de classification}
\begin{tabularx}{\textwidth}{|l|X|l|}
\hline
\rowcolor{lightblue}
\textbf{Métrique} & \textbf{Signification} & \textbf{Formule simplifiée} \\
\hline
Accuracy & Proportion de prédictions correctes & (TP + TN) / Total \\
\hline
Précision & Parmi les pannes prédites, combien sont réelles & TP / (TP + FP) \\
\hline
Rappel & Parmi les vraies pannes, combien sont détectées & TP / (TP + FN) \\
\hline
F1-Score & Équilibre entre précision et rappel & 2 $\times$ P $\times$ R / (P + R) \\
\hline
ROC AUC & Capacité discriminante globale du modèle & Aire sous la courbe ROC \\
\hline
\end{tabularx}
\end{table}

Dans le contexte de la maintenance prédictive, le \textbf{rappel} est la métrique
la plus critique : manquer une vraie panne (faux négatif) est bien plus coûteux
qu'une fausse alarme (faux positif). C'est pourquoi notre LSTM est configuré pour
maximiser le rappel, atteignant 100\% sur le dataset Backblaze.

% ─────────────────────────────────────────────────────────────
\section{Travaux connexes}
% ─────────────────────────────────────────────────────────────

Plusieurs travaux académiques ont abordé la prédiction de défaillance des disques
durs à partir des données SMART. \cite{lei2018machinery} présente une revue
complète des applications du ML à la maintenance prédictive industrielle. Les
approches basées sur les données Backblaze ont montré que les attributs SMART 5,
187, 197 et 198 sont les plus prédictifs de défaillance imminente, avec des
précisions allant de 85\% à 95\% selon les modèles \cite{pinheiro2007smart}.

Notre contribution se distingue par l'intégration d'un pipeline complet de bout
en bout : de la collecte automatisée des données jusqu'à l'interface
conversationnelle, en passant par deux modèles ML complémentaires (Random Forest
pour les métriques système agrégées et LSTM pour les séquences temporelles SMART),
le tout déployé en microservices Docker.

% ─────────────────────────────────────────────────────────────
\section{Conclusion}
% ─────────────────────────────────────────────────────────────

Cet état de l'art a présenté les fondements théoriques des technologies utilisées
dans notre projet. Le Random Forest et le LSTM constituent des approches
complémentaires pour la prédiction de pannes : le premier opère sur 63
caractéristiques statistiques agrégées sur 30 jours, le second capture les
dynamiques temporelles de dégradation des disques sur des séquences de 5 mesures
SMART. Les métriques ROUGE permettent d'évaluer objectivement la qualité du
chatbot sur un dataset de 20 questions de référence. Le chapitre suivant détaille
la conception architecturale du système.
